% from ega4-11-fr.tex
\oldpage[IV-3]{157}

Mais $(\mathcal F_1)_{x_1}=\mathcal F_x\otimes_{\mathcal O_x}\mathcal O_{x_1}$,
et $\mathcal O_{x_1}$ est un $\mathcal O_x$-module \emph{fidèlement plat}~;
donc (2.2.11, (iii)) $\mathcal F_x$ est un $\mathcal O_y$-module plat. Comme
$X'_1=X_1\times_{Y_1}Y'_1$, $\mathcal F'_1=\mathcal F_1\otimes_{Y_1}Y'_1$,
on est bien ramené à la situation de l'énoncé (11.6.1), avec $A$ et $A'$
remplacés par $A_1$ et $B_1$ respectivement.

\medskip
\noindent\textbf{IV)} \emph{Fin de la démonstration.} --- On est finalement
réduit à prouver (11.6.1) sous les hypothèses supplémentaires suivantes~:
\begin{enumerate}[label=(\roman*)]
  \item $A$ et $A'$ sont des anneaux locaux de $\mathbf Z$-algèbres de type fini
  (donc des anneaux \emph{excellents} (7.8.3))~;
  \item $A$ est intégralement clos~;
  \item le corps résiduel $k'$ de $A'$ est une extension finie et radicielle du
  corps résiduel $k$ de $A$.
\end{enumerate}
On sait alors ((7.8.3) et (2.3.8)) que sous les conditions (i) et (ii), si
$\mathfrak m$ et $\mathfrak m'$ sont les idéaux maximaux de $A$ et $A'$
respectivement, la topologie $\mathfrak m$-adique sur $A$ est induite par la
topologie $\mathfrak m'$-adique de $A'$. Le complété
$\widehat u:\widehat A\to\widehat{A'}$ est donc un homomorphisme injectif.
D'autre part, comme le morphisme
$\operatorname{Spec}(k')\to\operatorname{Spec}(k)$ est radiciel, il en est de
même du morphisme $f'^{-1}(y')\to f^{-1}(y)$ (\textbf I, 3.5.7), et il n'y a
donc qu'un seul point $x'$ de $X'$ dont les projections dans $X$ et $Y'$ soient
$x$ et $y'$ respectivement. On peut donc appliquer le résultat de (11.5.2).
C.Q.F.D.
\end{proof}

\begin{corollary}[11.6.2]
\label{IV.11.6.2-fr}
Soient $A$ un anneau local intègre unibranche, $Y=\operatorname{Spec}(A)$,
$f:X\to Y$ un morphisme localement de présentation finie, $\mathcal F$ un
$\mathcal O_X$-Module quasi-cohérent de présentation finie. Soient $A'$ un
anneau local, $A\to A'$ un homomorphisme local injectif~; on pose
$Y'=\operatorname{Spec}(A')$, $X'=X\times_Y Y'$, $f'=f_{(Y')}$,
$\mathcal F'=\mathcal F\otimes_Y Y'$. Soit $x$ un point de $X$ dont la
projection $f(x)=y$ est le point fermé de $Y$~; on suppose que $\mathcal F'$ est
$f'$-plat en tous les points $x'$ de $X'$ dont les projections dans $X$ et $Y'$
sont respectivement $x$ et le point fermé $y'$ de $Y'$. Alors $\mathcal F$ est
$f$-plat au point $x$.
\end{corollary}

\begin{proof}
On peut en effet reprendre la partie I) de la démonstration de (11.6.1), qui
prouve (avec les mêmes notations) que si $\mathcal F_1$ est $f_1$-plat en tous
les points $x_1$ de $X_1$ dont les projections respectives dans $X$ et $Y_1$
sont $x$ et $y_1$, alors $\mathcal F$ est $f$-plat au point $x$~; on est ainsi
ramené au cas où $A$ est intégralement clos, donc géométriquement unibranche, et
la conclusion résulte alors de (11.6.1).
\end{proof}

\subsection{Contre-exemples}
\label{subsection:IV.11.7-fr}

\begin{env}[11.7.1]
\label{IV.11.7.1-fr}
Considérons d'abord le cas où $A$ est un anneau local \emph{artinien}, et où les
hypothèses de (11.4.11) sont remplies \emph{sauf} la condition (ii) concernant
le corps résiduel $k$ de $A$. Nous allons voir que la conclusion de (11.4.11)
peut alors être en défaut. Soit $k$ un corps admettant une extension
\emph{galoisienne} $k'$ de degré $[k':k]>1$, et désignons par $\Gamma$ le groupe
de Galois de $k'$. Soit $A$ une $k$-algèbre ayant une base de 3 éléments
$1,a,b$ avec la table de multiplication $a^2=ab=ba=b^2=0$, de sorte que $A$ est
un anneau local artinien dont l'idéal maximal $\mathfrak m=ka+kb$ est de carré
nul. Soit $A'=A\otimes_k k'$, qui est une $k'$-algèbre de base $1,a,b$, anneau
local artinien d'idéal maximal
$\mathfrak m'=k'a+k'b=\mathfrak m A'$, de carré nul~; $A$ s'identifie
canoniquement à un sous-anneau de $A'$. Soit $\mathfrak J$ le sous-$k'$-espace
vectoriel de $\mathfrak m'$ engendré par $a+\gamma b$, où $\gamma\in k'$
n'appartient pas à $k$~; il est clair que $\mathfrak J$ est un idéal de $A'$.
Posons $B=A'/\mathfrak J$~; c'est un anneau artinien qui est un $A$-module
\emph{non plat}~; sinon (Bourbaki, \emph{Alg. comm.}, chap.~II, \S~3,
n$^\circ$~2, cor.~2 de la prop.~5), $B$ serait un $A$-module \emph{libre}~;
comme $A'$ est lui aussi un $A$-module libre, et que l'homomorphisme canonique
$A'/\mathfrak m A'\to B/\mathfrak m B$ est bijectif, l'homomorphisme canonique
$A'\to B=A'/\mathfrak J$ serait lui
% IV3_B23_P157_END

